\section{System Model and Methodology}

\subsection{Network and Trust Model}

We consider an enterprise validium architecture in which a single
organization operates an off-chain node that: (1)~receives state-changing
transactions from enterprise applications, (2)~maintains a Sparse Merkle
Tree of key-value state, (3)~aggregates transactions into fixed-size
batches, (4)~generates a Groth16 zero-knowledge proof of correct batch
execution, and (5)~submits the proof and new state root to an
Avalanche Layer~1 blockchain for on-chain verification.

The threat model assumes a trusted operator (the enterprise itself) but
requires that the L1 contract can independently verify that the claimed
state root is the result of applying valid state transitions to the
previous root. The ZK proof provides this guarantee without revealing
the transaction data.

\subsection{Circuit Design}

\subsubsection{MerklePathVerifier Template}

The \texttt{MerklePathVerifier(depth)} template computes a Merkle root
from a leaf hash, an array of sibling hashes, and path direction bits.
At each level $i$, it uses a \texttt{Mux1} selector to order the current
intermediate hash and the sibling according to the path bit, then hashes
the pair with \texttt{Poseidon(2)}. The output is the computed root.

\begin{equation}
h_{i+1} = \text{Poseidon}\!\left(
  \text{Mux}(h_i, s_i, b_i),\;
  \text{Mux}(s_i, h_i, b_i)
\right)
\label{eq:merkle-level}
\end{equation}

where $h_0$ is the leaf hash, $s_i$ is the sibling at level $i$, and
$b_i \in \{0,1\}$ indicates whether the current node is the left or
right child.

\subsubsection{ChainedBatchStateTransition Template}

The main circuit template, \texttt{ChainedBatchStateTransition(depth,\allowbreak
batchSize)}, processes a batch of $n$ sequential state transitions.
For each transaction $i \in [0, n)$:

\begin{enumerate}
\item Compute old leaf hash: $\ell_{\text{old}} = \text{Poseidon}(k_i, v_i^{\text{old}})$
\item Compute new leaf hash: $\ell_{\text{new}} = \text{Poseidon}(k_i, v_i^{\text{new}})$
\item Verify old Merkle path: assert $\texttt{MerklePathVerifier}(\ell_{\text{old}}, \vec{s}_i, \vec{b}_i) = r_i$
\item Compute new root: $r_{i+1} = \texttt{MerklePathVerifier}(\ell_{\text{new}}, \vec{s}_i, \vec{b}_i)$
\end{enumerate}

The root chain is initialized with $r_0 = \texttt{prevStateRoot}$ and
must terminate at $r_n = \texttt{newStateRoot}$. Both endpoints are
public inputs, along with \texttt{batchNum} and \texttt{enterpriseId}.
All per-transaction inputs (keys, values, siblings, path bits) are
private.

The critical design decision is \emph{sibling sharing}: old and new
Merkle path verifications use the same sibling array $\vec{s}_i$
and path bits $\vec{b}_i$, since only the leaf value changes. This
ensures consistency---the new root is computed over the same tree
structure, differing only in the updated leaf.

\subsubsection{Per-Transaction Constraint Decomposition}

Each transaction instantiates:

\begin{itemize}
\item 2 Poseidon(2) leaf hashes: $2 \times 240 = 480$ constraints
\item 2 MerklePathVerifier($d$) instances: each requires $d$ levels of
  (Poseidon(2) + 2$\times$Mux1), yielding $2 \times 519d$ constraints
\item 1 IsEqual root check: $\sim$5 constraints
\end{itemize}

The theoretical per-transaction cost is therefore:
\begin{equation}
C_{\text{tx}}^{\text{theory}} = 2 \times 519d + 480 + 5 = 1{,}038d + 485
\label{eq:theory-cost}
\end{equation}

\subsection{Benchmark Methodology}

We compiled and benchmarked seven circuit configurations spanning three
tree depths ($d \in \{10, 20, 32\}$) and three batch sizes
($n \in \{4, 8, 16\}$), with the $d{=}32, n{=}16$ configuration
excluded due to prohibitive key generation time. For each configuration:

\begin{enumerate}
\item \textbf{Compilation}: Circom~2.2.3 compiles the circuit to R1CS
  and WASM, reporting the constraint count.
\item \textbf{Input generation}: A Node.js script constructs a valid
  Poseidon-based SMT, inserts random key-value pairs, performs $n$
  sequential updates, and records the witness (keys, old/new values,
  Merkle siblings, path bits, intermediate roots).
\item \textbf{Witness generation}: SnarkJS~0.7.6 generates the full
  witness from the input JSON using the WASM binary.
\item \textbf{Key generation}: A Groth16 trusted setup is performed
  using an appropriate Powers of Tau ceremony (power selected to
  exceed the circuit's constraint count).
\item \textbf{Proof generation}: SnarkJS generates the Groth16 proof.
\item \textbf{Verification}: SnarkJS verifies the proof against the
  verification key and public inputs.
\end{enumerate}

All benchmarks were executed on a commodity desktop running Windows~11,
Node.js~v22.13.1, Circom~2.2.3, and SnarkJS~0.7.6. Timing was measured
via \texttt{Date.now()} surrounding each SnarkJS API call. Each
configuration was run once; stochastic variance analysis is deferred
to Stage~2 of the experiment protocol.

\subsection{Extrapolation Method}

For configurations beyond our direct benchmark capability (batch sizes
32, 64, 128), we extrapolate using:

\begin{itemize}
\item \textbf{Constraint count}: The exact formula derived from
  measurements (Section~\ref{sec:results}).
\item \textbf{Proving time}: A linear model of 65~$\mu$s per constraint
  for SnarkJS on commodity hardware, derived from the mean across all
  seven benchmarks. For Rapidsnark, we apply the published 4--10$\times$
  speedup factor from~\cite{tachyon2024,zkmopro2024}.
\end{itemize}
